\documentclass{article}
\usepackage{tcee-repaso}
\tcetema[Contabilidad y convergencia]{3.A.45}{Evidencia empírica sobre el crecimiento económico y la distribución de la renta entre los factores de producción. Contabilidad del crecimiento, con especial referencia a la productividad total de los factores. Convergencia económica internacional.}
\tcefecha{03/11/2026}
\begin{document}
\cabecerarepaso
\begin{columnas}

% ==========================================================================
\bloque{Esquema del cante}
\esquemacante{%
\apartadocante{Introducción: «por sus frutos los conoceréis»}{2,5}
\apartadocante{I. Hechos del crecimiento y distribución de la renta}{7}
\subcante{1. Algunos hechos empíricos}{2}
\subcante{2. \autor{Kaldor} y los nuevos hechos}{2,5}
\subcante{3. Distribución de la renta entre los factores}{2,5}
\apartadocante{II. Contabilidad del crecimiento y PTF}{9,5}
\subcante{1. Enfoques clásicos: primal (3) y dual (1)}{4}
\subcante{2. Enfoque econométrico}{0,5}
\subcante{3. Contabilidad del desarrollo}{2}
\subcante{4. Problemas}{1}
\subcante{5. Evidencia sobre la PTF}{2}
\apartadocante{III. Convergencia y factores del crecimiento}{10}
\subcante{1. Convergencia $\beta$ y $\sigma$}{2}
\subcante{2. Evidencia}{3,5}
\subcante{3. Causas próximas}{2}
\subcante{4. Causas últimas: las instituciones}{2,5}
\apartadocante{Conclusión}{1}}
\emph{Pizarra}: izquierda, residuo de \autor{Solow} y regresión de convergencia (hasta el final); centro, el gráfico de cada bloque; derecha, \autor{Kaldor} con $\checkmark$/$\times$, \autor{Hall}-\autor{Jones} y causas últimas. Corea del Sur vs. Nigeria en cada bloque. Pegatinas: «las apariencias engañan», «unos cardan la lana\ldots», «no se ganó Zamora en una hora», «Dios los cría\ldots», «árbol que nace torcido\ldots».

% ==========================================================================
\bloque[2,5]{Introducción}
\begin{itemize}
\item \textbf{Enganche}: «Por sus frutos los conoceréis». Nogales, partida por una valla: misma gente, clima e historia hasta 1853; Arizona, renta mediana $\approx$30.000\,\$ (hoy $>$43.000); Sonora, 1/3, sin escuela, mortalidad infantil alta (\autor{Acemoglu}-\autor{Robinson}, 2012). Las separan las \emph{instituciones}. ¿Por qué unos países son más ricos y crecen y otros no?
\item \textbf{Relevancia}: \emph{interés compuesto} (\autor{Lucas}, 1988: difícil pensar en otra cosa); \emph{pobreza} $<$3\,\$/día 43\,\% (1990) $\to$ 10\,\% (2024), India 52 $\to$ 5\,\%; la \emph{evidencia ordena la teoría} (Nobel 2024 a \autor{Acemoglu}, \autor{Johnson}, \autor{Robinson}); \emph{reparto}: ↓ del trabajo desde los 80.
\item \textbf{Contextualización}: \concepto{crecimiento}: ↑ sostenido de la producción pc \vertema{3.A.43} \vertema{3.A.44}. \concepto{Distribución funcional}: entre factores ($\neq$ personal). \concepto{Contabilidad del crecimiento}: más factores vs. usarlos mejor (\concepto{PTF}). \concepto{Convergencia}: pobres crecen más; \emph{absoluta} o \emph{condicional} (mismos determinantes del EE). \concepto{Causas próximas} ($K$, $H$, tecnología) vs. \concepto{últimas} (instituciones, geografía, cultura, suerte). Historia: \autor{Smith}, \autor{Ricardo}, \autor{Malthus}; \autor{Solow}-\autor{Swan} (1956); \autor{Solow} (1957); \autor{Kaldor} (1961); \autor{Jorgenson}-\autor{Griliches} (1967); \autor{Summers}-\autor{Heston} (1988), \autor{Baumol} (1986); \autor{Barro} (1991), MRW (1992), \autor{Hall}-\autor{Jones} (1999); AJR (2001); \autor{Karabarbounis}-\autor{Neiman} y \autor{Piketty} (2014).
\item \textbf{Problemática}: ¿se mantienen los hechos de \autor{Kaldor}, y el reparto estable? ¿Cuánto es acumulación y cuánto PTF, y qué mide esta? ¿Se igualan las rentas pc? Si no, ¿por qué?
\item \textbf{Estructura}: hechos $\to$ medida $\to$ causas: \hyperref[rep:I]{I}, hechos y reparto; \hyperref[rep:II]{II}, contabilidad y PTF; \hyperref[rep:III]{III}, convergencia, de las causas próximas a las últimas.
\end{itemize}

% ==========================================================================
\bloque[7]{I. Hechos del crecimiento y distribución de la renta\label{rep:I}}
\begin{ideaclave}
La frontera crece a un ritmo notablemente estable; los demás difieren muchísimo en nivel y ritmo. Casi todos los hechos de \autor{Kaldor} se mantienen salvo uno: la participación del trabajo cae desde los 80, sobre todo en EE.\,UU.
\end{ideaclave}

\sub{I.1. Algunos hechos empíricos\label{rep:I.1}}
\begin{itemize}
\item \emph{Frontera}: EE.\,UU., PIB pc ×12 en 1870-2022, 1,6\,\% anual (Maddison). Con 1 punto menos, ×$<$3 (como Indonesia, Perú o Egipto); con 1 más, ×56. Casi una recta en escala log.
\end{itemize}
\graficorepaso[0.9\linewidth]{pib-pc-largo}{PIB pc 1870-2022 (escala log.): EE.\,UU. casi recta; Japón, España y Corea se acercan; India despega tarde}
\begin{itemize}
\item \emph{Niveles}: EE.\,UU./Etiopía 36 (1960); EE.\,UU./Sudán del Sur 117 (2023). \emph{Ritmos}: 111 países, 1960-2023, media 1,9\,\%, desv. típica 1,3 puntos: China 6,1\,\%, Corea 5,4\,\%, R.\,D. Congo $-1{,}5$\,\%. \concepto{Milagros} ($>$3\,\% sostenido) y \concepto{desastres} (\autor{Panagariya}, 2004).
\item \emph{Corea del Sur vs. Nigeria}: 1.400 y 2.200\,\$ en 1960; 5,4 y 0,5\,\% anual; en 2023, 52.400 (más que España, 46.700) y 4.600: 11 veces. ¿Por qué?
\end{itemize}

\sub{I.2. \autor{Kaldor} y los nuevos hechos\label{rep:I.2}}
\begin{itemize}
\item \concepto{Hechos estilizados} (\autor{Kaldor}, 1961), a LP: (1) $Y/L$ crece a tasa $\approx$ constante; (2) $K/L$ también; (3) rendimiento real del capital $\approx$ constante; (4) $K/Y$ $\approx$ constante; (5) participaciones estables; (6) tasas muy distintas entre países. (1)+(4) $\Rightarrow$ (2); (3)+(4) $\Rightarrow$ (5), pues $\alpha=c_KK/Y$, $c_K=r+\delta$.
\item \concepto{Senda de crecimiento equilibrado} (todo a tasas constantes; en \autor{Solow}, el EE) \vertema{3.A.43}: reproduce 1-5 ($\gamma_{Y/L}=\gamma_{K/L}=g$, $r^*$, $k^*/f(k^*)$, $\alpha=k^*f'(k^*)/f(k^*)$); el 6.º, solo en la transición.
\item \emph{Datos} 1960-2023 (PWT): EE.\,UU. cumple ($\gamma_{Y/L}$ 2\,\%, $K/Y$ 3,8 $\to$ 3,2, rentabilidad $\approx$8\,\%); España, $K/Y$ 4,0 $\to$ 6,2 y rentabilidad 9,9 $\to$ 6,7\,\%: convergía. Con vivienda, riqueza/renta 2-3 $\to$ 4-6 (\autor{Piketty}-\autor{Zucman}, 2014). El 5.º, el que más cambia.
\item \concepto{Nuevos hechos} (\autor{Jones}-\autor{Romer}, 2010): mercados más amplios; crecimiento acelerado; varianza ↑ con la distancia a la frontera; enormes diferencias de renta y PTF; $h$ ↑ en todo el mundo; salarios relativos estables $\Rightarrow$ ideas \vertema{3.A.44}, PTF e instituciones.
\end{itemize}

\sub{I.3. Distribución de la renta entre los factores\label{rep:I.3}}
\begin{itemize}
\item Son el 5.º hecho y las ponderaciones de la contabilidad. \emph{Neoclásica}: RCE y pago por PMg; \autor{Euler} \vertema{3.A.11}:
\begin{formula}
\centering $Y=c_KK+wL,\qquad \alpha\equiv\dfrac{c_KK}{Y},\quad 1-\alpha\equiv\dfrac{wL}{Y}$
\end{formula}
\item[] \emph{Cobb-Douglas}: $\alpha$ fija ($\sigma=1$: lo que gana en cantidad lo pierde en precio; \autor{Cobb}-\autor{Douglas}, 1928). \emph{CES}, $\rho\equiv(\sigma-1)/\sigma$:
\begin{formula}
\centering $\alpha_t=a_K\left(K_t/Y_t\right)^{\rho}$:\ ↑ con $K/Y$ si $\sigma>1$, ↓ si $\sigma<1$
\end{formula}
\item[] Caída duradera: $\sigma\neq1$, progreso sesgado (automatización) o márgenes (renta de monopolio).
\item \autor{Piketty} (2014): $\alpha=r\,v$ y $v\to s/\gamma_Y$; $r>g$ concentra la riqueza. Críticas: exige $\sigma>1$ y es vivienda (\autor{Rognlie}, 2015); ahorro neto (\autor{Krusell}-\autor{Smith}, 2015); sin instituciones; $r>g$ casi siempre (\autor{Mankiw}, 2015) \vertema{3.A.24}.
\item \emph{Medida}: autónomos (España 2025: 48,4\,\% sin ajustar, 56,6\,\% con la \concepto{participación salarial ajustada}; \autor{Gollin}, 2002: 65-80\,\% $\Rightarrow$ $\alpha\approx1/3$); vivienda y depreciación.
\item \emph{Evidencia}: estable hasta los 80 (capital 0,3-0,5); $\approx-5$ puntos en el mundo, dentro de cada sector (\autor{Karabarbounis}-\autor{Neiman}, 2014). EE.\,UU. 63 $\to$ 54\,\%; España 67,5\,\% (1976) $\to$ 55\,\% (2007) $\to$ 53-58\,\%; zona del euro 58,4 $\to$ 54,2 $\to$ 56,9\,\%. Con series homogéneas, solo EE.\,UU. y Canadá (\autor{Gutiérrez}-\autor{Piton}, 2020).
\end{itemize}
\graficorepaso[0.9\linewidth]{participacion-trabajo}{Participación salarial ajustada, 1960-2025 (\% del PIB; AMECO)}
\begin{itemize}
\item \emph{Explicaciones}: \emph{precio del capital} ($\sigma\approx1{,}25$ explica la mitad; pero $\sigma<1$ casi siempre); \emph{globalización} (y 1/3 de EE.\,UU. es imputación de autónomos: \autor{Elsby}-\autor{Hobijn}-\autor{Şahin}, 2013); \emph{superestrellas} (\autor{Autor} et al., 2020); \emph{márgenes} 21 $\to$ 61\,\% (\autor{De Loecker}-\autor{Eeckhout}-\autor{Unger}, 2020): caen \emph{las 2} (\autor{Barkai}, 2020) \vertema{3.A.15}; \emph{sindicatos}. Ninguna lo explica todo.
\end{itemize}

% ==========================================================================
\bloque[9,5]{II. Contabilidad del crecimiento y PTF\label{rep:II}}
\begin{ideaclave}
La acumulación explica buena parte del crecimiento en el tiempo, pero no la mayor parte de las diferencias de renta entre países, que son sobre todo productividad. La PTF es un residuo: mide todo lo que no medimos.
\end{ideaclave}

\sub{II.1. Enfoques clásicos: primal y dual\label{rep:II.1}}
\begin{itemize}
\item \textbf{Primal} (\autor{Solow}, 1957): $Y=A\,F(K,L)$ \concepto{neutral de \autor{Hicks}} (no cambia la RMST), RCE, pago por PMg $\Rightarrow$ elasticidades = participaciones. \concepto{Residuo de \autor{Solow}}:
\funcionclave{\gamma_A=\gamma_Y-\alpha\,\gamma_K-(1-\alpha)\,\gamma_L}
\item[] Por hora: $\gamma_{Y/L}=\gamma_A+\alpha\,\gamma_{K/L}$; datos anuales, $\bar\alpha_t$ media de 2 años (\concepto{\autor{Törnqvist}}). \autor{Harrod} $Y=K^\alpha(AL)^{1-\alpha}$ = \autor{Hicks} con PTF $A^{1-\alpha}$.
\item[] EE.\,UU. 1909-1949: $Y/L$ ×2; capital 1/8, residuo 7/8. Residuo = técnica + organización + asignación + calidad no medida + errores. Choque de los ciclos reales (\autor{Kydland}-\autor{Prescott}, 1982) \vertema{3.A.42}. Críticas: competencia perfecta; progreso incorporado (\autor{Kaldor}, 1957); «\concepto{medida de nuestra ignorancia}» (\autor{Abramovitz}, 1956); no es causalidad.
\item \textbf{Dual} (\autor{Hsieh}, 2002): con \emph{precios}; solo $Y=c_KK+wL$:
\begin{formula}
\centering $\underbrace{\gamma_Y-\alpha\gamma_K-(1-\alpha)\gamma_L}_{\text{primal}}=\underbrace{\alpha\,\gamma_{c_K}+(1-\alpha)\,\gamma_w}_{\text{dual}}=\gamma_A$
\end{formula}
\item[] Sin progreso, sobre la \concepto{frontera de precios de los factores}: ↑$w$ $\Rightarrow$ ↓$c_K$. Identidad: si difieren, fallan los datos.
\item[] \emph{Tigres}: \autor{Young} (1995): acumulación; Singapur, PTF 0,2\,\% (1966-1990); \autor{Krugman} (1994), como la URSS: «las apariencias engañan». \autor{Hsieh}: el rendimiento del capital apenas cae, PTF dual +2 puntos (inversión inflada). \autor{Young} (1998): dual $\approx$ primal.
\end{itemize}

\sub{II.2. Enfoque econométrico\label{rep:II.2}}
\begin{itemize}
\item $\gamma_{Y,t}=\beta_0+\beta_1\gamma_{K,t}+\beta_2\gamma_{L,t}+u_t$, $\beta_0$ = PTF media; no supone pago por PMg. Pero \emph{simultaneidad} (\autor{Marschak}-\autor{Andrews}, 1944); residuo no exógeno (\autor{Hall}, 1988); elasticidades que cambian.
\end{itemize}

\sub{II.3. Contabilidad del desarrollo\label{rep:II.3}}
\begin{itemize}
\item Niveles entre países. $Y=K^\alpha(AhL)^{1-\alpha}$, $h=e^{\phi(S)}$ (\autor{Mincer}) \vertema{3.A.25}:
\funcionclave{\frac{Y}{L}=\left(\frac{K}{Y}\right)^{\frac{\alpha}{1-\alpha}}h\,A}
\item[] $K/Y$, no $K/L$: el capital acumulado \emph{por} ser productivo va a la PTF (\autor{Klenow}-\autor{Rodríguez-Clare}, 1997).
\item \textbf{MRW} (1992), 98 países, 1960-1985, $g+\delta=0{,}05$:
\begin{formula}
\centering $\ln\dfrac{Y_i}{L_i}=a+\dfrac{\alpha\ln s_{K}+\beta\ln s_{H}-(\alpha+\beta)\ln(n+g+\delta)}{1-\alpha-\beta}+\varepsilon_i$
\end{formula}
\item[] Sin $H$: $R^2=0{,}59$, $\alpha=0{,}60$ (el doble); con $H$: $R^2=0{,}78$, $\alpha\simeq0{,}31$, $\beta\simeq0{,}28$; OCDE, pobre. Pero la matrícula exagera $H$ y supone $A$ independiente de $s_K$.
\item \textbf{\autor{Hall}-\autor{Jones}} (1999), calibran, 127 países, 1988, $\alpha=1/3$: 5 más ricos/5 más pobres, $31{,}7=1{,}8\,(K)\times2{,}2\,(h)\times8{,}3\,(A)$; correlación de logaritmos de $A$ e $Y/L$, 0,89. Causa: \concepto{infraestructura social}. 2023: $47{,}1=1{,}07\times2{,}29\times19{,}1$; Corea/Nigeria, $8{,}1=0{,}7\times1{,}8\times6{,}2$.
\end{itemize}
\graficorepaso[0.7\linewidth]{ptf-renta}{PTF y producción por trabajador, 2023, relativas a EE.\,UU. (escala log.): casi toda la diferencia es PTF}
\begin{itemize}
\item \textbf{\autor{Caselli}} (2005): \concepto{éxito} $=\mathrm{var}[\ln(Y/L)_{KH}]/\mathrm{var}[\ln(Y/L)]$: 0,39 (1996), 0,30 (2023). MRW: «casi todo acumulación»; desarrollo: «casi todo productividad».
\end{itemize}

\sub{II.4. Problemas\label{rep:II.4}}
\begin{itemize}
\item \emph{Factores}: \concepto{inventario permanente}, $K_{t+1}=(1-\delta)K_t+I_t$ (stock $\neq$ servicios: \concepto{índice de servicios}); horas y composición. Calidad no medida $\Rightarrow$ PTF: \autor{Jorgenson}-\autor{Griliches} (1967), residuo 1,6 $\to$ 0,1\,\%; debate con \autor{Denison}. Servicios y sector público.
\item \emph{Bruta o VAB}: la bruta depende de la integración vertical: VAB o \concepto{pesos de \autor{Domar}}.
\item \emph{Contabilidad $\neq$ causalidad}: en el EE de \autor{Solow}, da $\alpha g$ al capital, pero lo causa la tecnología: «unos cardan la lana y otros llevan la fama». Con $K/Y$:
\begin{formula}
\centering $\gamma_{Y/L}=\dfrac{\gamma_A}{1-\alpha}+\dfrac{\alpha}{1-\alpha}\,\gamma_{K/Y}$
\end{formula}
\end{itemize}

\sub{II.5. Evidencia sobre la PTF\label{rep:II.5}}
\begin{itemize}
\item \emph{Mundo} 1960-2000, 84 países: $Y/L$ 2,3\,\% = capital 1 + educación 0,3 + PTF 0,9 (\autor{Bosworth}-\autor{Collins}, 2003).
\item \emph{España} 1996-2007: 3,6\,\% = capital 1,8 + horas 2,0 + $h$ 0,4 + PTF $-0{,}5$; en 1971-1985, PTF 2 de 2,8 puntos \vertema{4.A.27}. EE.\,UU. 1996-2007, PTF 1,2 (TIC): la diferencia con Europa.
\end{itemize}
\graficorepaso[\linewidth]{contabilidad-espana}{Contabilidad del crecimiento (puntos, media anual): España 1996-2007, PTF $-0{,}5$}
\begin{itemize}
\item \emph{Desaceleración} (EE.\,UU.): 2,1\,\% (1948-1973), 0,5\,\% (1974-1995), 1,8\,\% (1996-2004), 0,5\,\% (2005-2019); \autor{Denison} (1985). Causas: petróleo, no; medición, no (\autor{Syverson}, 2017); servicios (\autor{Baumol}); ideas más difíciles (\autor{Bloom} et al., 2020; \autor{Gordon}, 2016); mala asignación (España); \concepto{curva en J} de la IA (\autor{Brynjolfsson}-\autor{Rock}-\autor{Syverson}, 2021).
\item PTF $\neq$ productividad aparente $Y/L$ (sube también al acumular; en España, en las recesiones, por composición).
\end{itemize}

% ==========================================================================
\bloque[10]{III. Convergencia y factores del crecimiento\label{rep:III}}
\begin{ideaclave}
Los pobres no alcanzaron a los ricos en la 2.ª mitad del s.\,XX; sí las economías parecidas, despacio ($\approx$2\,\% anual), y desde 2000 asoma algo en el mundo. Detrás, las causas próximas y, en último término, las instituciones.
\end{ideaclave}

\sub{III.1. Convergencia $\beta$ y $\sigma$\label{rep:III.1}}
\begin{itemize}
\item \concepto{Convergencia} real ($\neq$ nominal \vertema{3.B.40}): predicción de \autor{Solow}; en $AK$, no \vertema{3.A.44}.
\item \concepto{Convergencia $\beta$} (\autor{Barro}-\autor{Sala-i-Martín}, 1991):
\funcionclave{\frac1T\ln\frac{(Y/L)_{i,t}}{(Y/L)_{i,t-T}}=\beta_0+\beta_1\ln\Big(\frac YL\Big)_{i,t-T}+\vect{x}_i\cdot\vect{\psi}+u_{i,t}}
\item[] \concepto{Absoluta}: $\beta_1<0$ sin controles; \concepto{condicional}: controlando $\vect{x}_i$ (cada una a \emph{su} EE). $\beta_1=-(1-e^{-\lambda T})/T$, $\lambda=(1-\alpha)(n+g+\delta)$: con $\alpha=1/3$, $\lambda\simeq4$\,\% (17 años de semivida); con 2/3, $\simeq2$\,\% (35).
\item \concepto{Convergencia $\sigma$}: $\sigma_t<\sigma_{t-T}$. $\sigma\Rightarrow\beta$, $\beta\not\Rightarrow\sigma$: adelantamiento o perturbaciones:
\begin{formula}
\centering $\sigma_t^2=(1-b)^2\sigma_{t-1}^2+\sigma_u^2\ \to\ \sigma^{*2}=\dfrac{\sigma_u^2}{1-(1-b)^2}$
\end{formula}
\end{itemize}
\graficorepaso[0.9\linewidth]{beta-sigma}{(a) $\beta$ y $\sigma$; (b) ninguna; (c) $\beta$ sin $\sigma$: adelantamiento}
\begin{itemize}
\item \concepto{Clubes de convergencia} (trampas) \vertema{3.B.28}. \concepto{Falacia de \autor{Galton}}: regresión a la media; mirar toda la distribución (\autor{Quah}, 1993).
\item Predicen convergencia: neoclásica, Sobelow, \autor{Jones}; no: $AK$, \autor{Romer}; kaldorianos y nueva geografía, posible divergencia.
\end{itemize}

\sub{III.2. Evidencia\label{rep:III.2}}
\begin{itemize}
\item \autor{Baumol} (1986): 16 industrializados, 1870-1979, convergen; \autor{De Long} (1988): \emph{sesgo de selección} (con los ricos de 1870, desaparece). MRW: 98 países, coeficiente $+$; \autor{Pritchett} (1997): «divergencia a lo grande», ricos/pobres 8,7 $\to$ 45. Se leyó a favor de los endógenos.
\end{itemize}
\graficorepaso[\linewidth]{convergencia-periodos}{Convergencia $\beta$: 1960-2000 (todos $+0{,}25$; OCDE $-1{,}03$) y 2000-2023 (todos $-0{,}47$; OCDE $-1{,}49$)}
\begin{itemize}
\item \emph{Condicional} (\autor{Barro}-\autor{Sala-i-Martín}, 1992; MRW): con $s_K$ y $n$, $-0{,}14$; con escolarización, $-0{,}29$ (1,4\,\%); OCDE, $\approx$2\,\%. \concepto{Sesgo de variable omitida}: los pobres ahorran menos.
\item \emph{Regiones}: EE.\,UU., Japón, Europa, $\approx$2\,\% anual (media distancia en 35 años): «no se ganó Zamora en una hora»; exige capital amplio 0,7-0,8. Estados 1929-2025, 1,4\,\% \vertema{4.A.17}.
\item \emph{$\sigma$}: mundo 0,96 $\to$ 1,34 $\to$ 1,27 (no hay); OCDE 0,61 $\to$ 0,44; UE-27 0,55 $\to$ 0,33.
\item \concepto{Picos gemelos} (\autor{Quah}, 1996; \autor{Durlauf}-\autor{Johnson}, 1995): «Dios los cría y ellos se juntan».
\item \emph{Desde 2000}: absoluta (\autor{Kremer}-\autor{Willis}-\autor{You}, 2022); \autor{Acemoglu}-\autor{Molina} (2022): con efectos fijos, no. Solo 0,5\,\% anual: frágil. Manufacturas, incondicional (\autor{Rodrik}, 2013).
\end{itemize}

\sub{III.3. Causas próximas\label{rep:III.3}}
\begin{itemize}
\item \autor{Barro} (1991), 98 países: convergencia dado $h$ inicial; + capital humano; $-$ consumo público; $-$ inestabilidad política; $-$ distorsiones; ficticias de África y América Latina $<0$. Inflación alta, $-$ (1995).
\item Inversión; finanzas (\autor{King}-\autor{Levine}, 1993); comercio instrumentado con geografía (\autor{Frankel}-\autor{Romer}, 1999) \vertema{3.B.9}. Corea/Nigeria: $h$ 3,9 vs. 2,1; ¿por qué?
\item \emph{Robustez}: \autor{Levine}-\autor{Renelt} (1992), \concepto{cotas extremas} (\autor{Leamer}): solo inversión; \autor{Sala-i-Martín} (1997), millones de regresiones; bayesiano (2004): 18 de 67.
\end{itemize}

\sub{III.4. Causas últimas: las instituciones\label{rep:III.4}}
\begin{itemize}
\item 4 causas (\autor{Acemoglu}, 2009): suerte (equilibrios múltiples), geografía (\autor{Sachs}; \autor{Diamond}), cultura (\autor{Weber}; \autor{Tabellini}), instituciones. Contra la geografía: \concepto{inversión de la fortuna} (AJR, 2002). \concepto{Instituciones}: «las reglas del juego» (\autor{North}, 1990).
\item \textbf{AJR} (2001), \concepto{instrumento}: mortalidad de colonos $\Rightarrow$ \concepto{instituciones extractivas} $\Rightarrow$ persisten $\Rightarrow$ renta; \emph{exclusión}: «árbol que nace torcido\ldots». 64 colonias:
\begin{formula}
\centering $R_i=\zeta+\beta_1\ln M_i+v_i;\quad \ln(Y/L)_i=\mu+\beta_2\hat R_i+\varepsilon_i$
\end{formula}
\item[] $\hat\beta_2=0{,}94>$ MCO 0,52. Nigeria con instituciones de Chile, renta ×7. Críticas: \autor{Albouy} (2012), \autor{Glaeser} et al. (2004). 2 Coreas: el árbol se endereza.
\item «Las instituciones mandan» (\autor{Rodrik}-\autor{Subramanian}-\autor{Trebbi}, 2004); democracia, +20\,\% en 25 años (\autor{Acemoglu} et al., 2019); persisten por \emph{conflicto social} (AJR, 2005).
\end{itemize}

% ==========================================================================
\bloque[1]{Conclusión}
\begin{itemize}
\item \textbf{Recapitulación}: «En conclusión,» \autor{Kaldor} se cumple salvo el reparto; acumular explica el crecimiento en el tiempo, la PTF las diferencias entre países, y se ha frenado; sin convergencia absoluta en el s.\,XX, sí condicional ($\approx$2\,\%); detrás, las instituciones.
\item \textbf{Relevancia}: agenda de productividad (competencia, asignación, innovación, instituciones): informe \autor{Draghi} (2024); España, PTF casi plana desde los 90.
\item \textbf{Extensiones}: \vertema{3.A.43} \vertema{3.A.44} \vertema{3.B.40} \vertema{4.A.27} \vertema{3.B.28}.
\item \textbf{Opinión}: de «cuánto capital» a «cuánta productividad» y a «qué instituciones»; impulsos que se agotan (\autor{Jones}, 2023).
\item \textbf{Idea final}: Nogales, misma tierra y gente, frutos distintos según las reglas del juego: «por sus frutos los conoceréis».
\end{itemize}

% ==========================================================================
\bloque{Autores y fechas}
\hito{1928}{\autor{Cobb}-\autor{Douglas}: participaciones constantes}
\hito{1956}{\autor{Solow}, \autor{Swan}; \autor{Abramovitz}: «medida de nuestra ignorancia»}
\hito{1957}{\autor{Solow}: residuo (7/8)}
\hito{1961}{\autor{Kaldor}: hechos estilizados}
\hito{1967}{\autor{Jorgenson}-\autor{Griliches}: calidad de los factores}
\hito{1986}{\autor{Baumol}; 1988, \autor{De Long}: sesgo de selección}
\hito{1991}{\autor{Barro}: regresiones; \autor{Barro}-\autor{Sala-i-Martín}: $\beta$ y $\sigma$}
\hito{1992}{MRW; \autor{Levine}-\autor{Renelt}}
\hito{1994}{\autor{Krugman}: el mito del milagro asiático}
\hito{1995}{\autor{Young}: tigres por acumulación}
\hito{1996}{\autor{Quah}: picos gemelos}
\hito{1997}{\autor{Pritchett}: divergencia a lo grande}
\hito{1999}{\autor{Hall}-\autor{Jones}: infraestructura social}
\hito{2001}{AJR: mortalidad de colonos}
\hito{2002}{\autor{Hsieh}: dual}
\hito{2005}{\autor{Caselli}: éxito de los factores}
\hito{2010}{\autor{Jones}-\autor{Romer}: nuevos hechos}
\hito{2014}{\autor{Karabarbounis}-\autor{Neiman}; \autor{Piketty}}
\hito{2020}{\autor{Bloom} et al.: ideas más difíciles; \autor{De Loecker} et al.: márgenes}
\hito{2022}{\autor{Kremer}-\autor{Willis}-\autor{You}; \autor{Acemoglu}-\autor{Molina}}
\hito{2024}{Nobel a \autor{Acemoglu}, \autor{Johnson} y \autor{Robinson}}

% ==========================================================================
\bloque{Datos clave}
\begin{itemize}
\item PIB pc EE.\,UU. 1870-2022: \dato{×12, 1,6\,\% anual}{Maddison, 2025}.
\item Pobreza $<$3\,\$/día: \dato{43\,\% (1990) $\to$ 10\,\% (2024)}{BM, 2026}.
\item Corea/Nigeria 2023: \dato{52.400 vs. 4.600\,\$}{PWT 11.0}; España \dato{46.700\,\$}{PWT}.
\item Participación salarial ajustada 2025: España \dato{56,6\,\%}{AMECO}; zona del euro \dato{56,9\,\%}{AMECO}; EE.\,UU. \dato{54\,\%}{AMECO}.
\item Remuneración de asalariados, España 2025: \dato{48,4\,\% del PIB}{Eurostat}.
\item Márgenes de las cotizadas de EE.\,UU.: \dato{21\,\% (1980) $\to$ 61\,\% (2016)}{\autor{De Loecker} et al., 2020}.
\item PTF España 1996-2007: \dato{$-0{,}5$ puntos}{AMECO, PWT}; EE.\,UU. \dato{0,9\,\% desde 1970}{FRBSF}.
\item Esfuerzo investigador \dato{×23}{\autor{Bloom} et al., 2020}; productividad de la investigación \dato{÷41, $-5$\,\%/año}{ídem}.
\item Desarrollo 1988: \dato{$31{,}7=1{,}8\times2{,}2\times8{,}3$}{\autor{Hall}-\autor{Jones}, 1999}.
\item Desarrollo 2023: \dato{$47{,}1=1{,}07\times2{,}29\times19{,}1$}{PWT}; éxito de los factores \dato{0,30}{PWT}.
\item $\sigma$ mundial: \dato{0,96 $\to$ 1,34 $\to$ 1,27}{PWT, 1960-2023}.
\end{itemize}

% ==========================================================================
\bloque{Preguntas del tribunal}
\begin{enumerate}
\item \textbf{¿Qué mide la PTF?} Un residuo: técnica, organización, asignación, calidad no medida y errores.
\item \textbf{¿Por qué primal y dual dan lo mismo?} Identidad $Y=c_KK+wL$; si difieren, fallan cantidades o precios (Singapur).
\item \textbf{¿Por qué $K/Y$ en \autor{Hall}-\autor{Jones}?} El capital acumulado por ser productivo se atribuye a la PTF.
\item \textbf{¿Implica la $\beta$ la $\sigma$?} No (adelantamiento, perturbaciones); $\sigma\Rightarrow\beta$.
\item \textbf{¿Por qué la regresión absoluta oculta la convergencia?} Variable omitida: los pobres ahorran menos y su población crece más.
\item \textbf{¿Por qué cae la participación del trabajo?} Precio del capital ($\sigma>1$, discutido), globalización, superestrellas, márgenes, sindicatos; clara solo en EE.\,UU.
\item \textbf{¿Es causal la contabilidad?} No: en el EE de \autor{Solow} el capital crece porque crece la PTF; corregir con $K/Y$.
\item \textbf{¿Por qué la convergencia es tan lenta?} $\approx$2\,\% frente al 4\,\% de \autor{Solow} con $\alpha=1/3$: exige un capital amplio (con $H$) de 0,7-0,8; los rendimientos decrecientes llegan despacio.
\item \textbf{¿Instituciones o geografía?} Las instituciones mandan (\autor{Rodrik} et al., 2004); la geografía actúa sobre todo a través de ellas (inversión de la fortuna), aunque la malaria pesa (\autor{Sachs}, 2003).
\item \textbf{¿Qué supuesto clave en AJR?} Exclusión: la mortalidad solo afecta a la renta vía instituciones; críticas de \autor{Albouy} y \autor{Glaeser}.
\end{enumerate}

\begin{trampa}
\begin{itemize}
\item \autor{Kaldor}: bastan 1, 3, 4 y 6 (1+4 $\Rightarrow$ 2; 3+4 $\Rightarrow$ 5); \autor{Solow} no reproduce el 6.º salvo en la transición.
\item Residuo: $\alpha$ = participación del \emph{capital}; \autor{Solow} (1957), capital 1/8, residuo 7/8.
\item $\sigma\Rightarrow\beta$, no al revés. Absoluta: $\beta_1<0$ \emph{sin} controles.
\item Velocidad empírica $\approx$2\,\%, la mitad del 4\,\% de \autor{Solow} con $\alpha=1/3$.
\item MRW sin $H$: $\alpha=0{,}60$; con $H$: 0,31 y $\beta$ 0,28.
\item Cobb-Douglas: participaciones constantes; CES: $\alpha$ ↑ con $K/Y$ solo si $\sigma>1$.
\item PTF de \autor{Hicks} $=A^{1-\alpha}$ si $A$ aumenta el trabajo: crece a $(1-\alpha)g$.
\item Dual (\autor{Hsieh}): sin forma funcional ni stock de capital; sí que los precios agoten el producto.
\item Contabilidad del desarrollo: «casi todo productividad» (\autor{Hall}-\autor{Jones}, \autor{Caselli}) vs. MRW, «casi todo acumulación».
\item \autor{De Long}: sesgo de \emph{selección}; la absoluta, de \emph{variable omitida}.
\item Remuneración de asalariados $\neq$ participación ajustada (autónomos).
\item AJR: VI (0,94) $>$ MCO (0,52); el instrumento explica el \emph{nivel}, no el crecimiento.
\end{itemize}
\end{trampa}

\end{columnas}
\end{document}
